\section{はじめに}

\subsection{何を知りたいのか}

スペースバルーンは、ヘリウムまたは水素を詰めたゴム気球で観測機器を吊り上げ、
上空 \SI{30}{\kilo\metre} 前後で気球が破裂したあとパラシュートで降りてくる。
打ち上げから回収までおよそ 2--3 時間、水平方向には数十から数百キロメートル流される。

このとき機体を運ぶのは風だけである。
上昇中も下降中も、機体は自分の力で進むわけではなく、
その高度の風にそのまま乗って流されていく。
したがって落下地点は、地上から破裂高度までの\textbf{全高度の風}の積み重ねで決まる。

打ち上げの計画を立てる段階で知りたいのは、次のような問いである。

\begin{itemize}
  \item 3 月の前半に十勝から上げたら、機体はふつうどちらへ何キロ流れるのか。
  \item 落下範囲のばらつきは、どの時期が小さいのか。
  \item 海に落ちる危険が小さいのはいつか。
  \item モンテカルロ計算で風に与えるべき「ばらつき」は、実際どれくらいなのか。
\end{itemize}

いずれも「明日の風」ではなく「その時期の風のふつうの姿」を問うている。
これを気候値（climatology）と呼ぶ。

\subsection{なぜ既存のデータでは答えられないのか}
\label{sec:intro-ceiling}

現在のシミュレータには風の入手経路が 2 つあるが、どちらも気候値を与えない。

\paragraph{Open-Meteo の気圧面 API（\texttt{backend/bfps/core/wind.py}）}
気圧面ごとの風を返してくれるが、扱える一番上の面が \SI{30}{\hecto\pascal}、
高度にして約 \SI{23.8}{\kilo\metre} までしかない。
実装はそれより上を「最上層と同じ風が続く」と仮定して引き伸ばしている。
破裂高度を \SI{30}{\kilo\metre} とすると、この空白は飛行時間の 17 \% にあたる。

この空白の影響の大きさは季節によってまったく違う。
本報告の結果を使って測ったところ、冬は水平移動の 1--5 \% にすぎないのに対し、
夏は北海道で 54 \%、和歌山で \textbf{111 \%}、愛媛では \textbf{140 \%} に達する（第 \ref{sec:ceiling} 節）。
しかも夏は、落下距離が年間で最も短くなるので、打ち上げに向いた季節である。
つまり\textbf{現在の実装は、打ち上げに最も向いた季節に最も大きな誤差を出す}。

\paragraph{Tawhiri / GFS（\texttt{frontend/src/api/tawhiri.ts}）}
こちらは全高度をカバーするが、数値予報の出力であり、
実装上も「12 時間前から 7 日後まで」に限られている。
来月の前半を問い合わせることはできない。

つまり、\textbf{過去の風を時期ごとにまとめたデータが手元に無い}。
本報告はそのデータを作る。

\subsection{本報告がしたこと}

気象庁の長期再解析 JRA-3Q から、

\begin{itemize}
  \item 期間 2016--2025 年（10 年）、
  \item 高度 地上から約 \SI{40}{\kilo\metre}（気圧 \SIrange{1000}{3}{\hecto\pascal} の 38 面）、
  \item 水平 \ang{0.375} 格子（約 \SI{40}{\kilo\metre} 間隔）、
  \item 時間 6 時間ごと（1 日 4 回）
\end{itemize}

のデータを取得し、北海道・和歌山・愛媛（愛南町）それぞれについて、
月を 2 つに割った半月ごと・高度ごと・格子点ごとに
風向・風速の統計をとった。

出力は NetCDF（本体）、Parquet（扱いやすい表形式）、要点だけの CSV、
そして本報告に載せた図である。

\subsection{この報告の読み方}

第 \ref{sec:data} 章でデータの素性を、第 \ref{sec:method} 章で計算方法を述べる。
方法の章がやや長いのは、手順に加えて
\textbf{なぜその方法でなければいけないのか}まで説明しているためである。
風向のような「円になっている量」は、ふつうの平均を取ると誤った答えが出る。
その理屈を飛ばすと、第 \ref{sec:hokkaido} 章以降の図の読み方を間違える。

第 \ref{sec:validation} 章は\textbf{この統計をどこまで信じてよいか}を扱う。
「月を 2 分割した半月」という単位は人間側の都合で決めたもので、
データがその粒度を支えているかは別の話である。
そこを 1 年抜き交差検証で確かめた。結果の図を見る前に読んでほしい。

結果は第 \ref{sec:hokkaido} 章（北海道）、第 \ref{sec:wakayama} 章（和歌山）、
第 \ref{sec:ehime} 章（愛媛）、第 \ref{sec:compare} 章（3 領域の比較）にある。
第 \ref{sec:limits} 章には、この結果で言えないことを書いた。
